% Part: many-valued-logic
% Chapter: three-valued-logics
% Section: lukasiewicz

\documentclass[../../../include/open-logic-section]{subfiles}

\begin{document}

\olfileid[id]{mvl}{thr}{luk}

\olsection{Logika \L ukasiewicz}

Salah satu usulan logika bernilai banyak pertama yang diterbitkan dan
dikembangkan secara terperinci diajukan oleh filsuf Polandia Jan
\L ukasiewicz pada tahun 1921. \L ukasiewicz dimotivasi oleh masalah
pertempuran laut Aristoteles: tampaknya, \emph{hari ini}, kalimat
``Besok akan terjadi pertempuran laut'' tidak benar maupun salah; nilai
kebenarannya belum ditentukan. \L ukasiewicz mengusulkan nilai kebenaran
ketiga bagi kalimat ``kontingen masa depan'' semacam itu.
\begin{quote}
Saya dapat mengandaikan tanpa kontradiksi bahwa keberadaan saya di Warsawa
pada suatu saat tertentu tahun depan, misalnya pada tengah hari tanggal
21~Desember, pada saat ini tidak ditentukan secara positif maupun negatif.
Jadi, mungkin saja saya berada di Warsawa pada waktu tersebut, tetapi hal itu
tidak niscaya. Berdasarkan pengandaian ini, proposisi ``Saya akan berada di
Warsawa pada tengah hari tanggal 21~Desember tahun depan'' pada saat ini tidak
benar maupun salah. Sebab, jika proposisi itu sekarang benar, keberadaan saya
kelak di Warsawa haruslah niscaya, yang bertentangan dengan pengandaian. Di
pihak lain, jika proposisi itu sekarang salah, keberadaan saya kelak di
Warsawa haruslah mustahil, yang juga bertentangan dengan pengandaian. Karena
itu, proposisi yang dibahas pada saat ini tidak benar maupun salah dan harus
mempunyai nilai ketiga, yang berbeda dari ``0'' atau kesalahan dan ``1'' atau
kebenaran. Nilai ini dapat kita lambangkan dengan~``$\frac{1}{2}$.'' Nilai
tersebut merepresentasikan ``yang mungkin'', dan menyertai ``yang benar'' dan
``yang salah'' sebagai nilai ketiga.
\end{quote}
Kita akan menggunakan $\Undef$ bagi nilai kebenaran ketiga
\L ukasiewicz.\footnote{Di sini \L ukasiewicz menggunakan ``mungkin'' dengan
arti yang tidak lazim dewasa ini, yakni mungkin tetapi tidak niscaya.}

Fungsi kebenaran bagi konektif $\lnot$, $\land$, dan~$\lor$ mudah ditentukan
berdasarkan penafsiran ini: negasi suatu kalimat kontingen masa depan juga
merupakan kalimat kontingen masa depan, sehingga
$\tf{\lnot}(\Undef)=\Undef$. Jika salah satu konjung dari suatu konjungsi
belum ditentukan dan konjung lainnya benar, konjungsi tersebut juga belum
ditentukan---bagaimanapun, bergantung pada hasil konjung kontingen masa depan,
konjungsi itu kelak mungkin benar dan mungkin pula salah. Jadi
\[
    \tf{\land}(\True,\Undef)=
\tf{\land}(\Undef,\True)=
\Undef.
\]
Namun, jika konjung lainnya salah, konjungsi tersebut tidak mungkin menjadi
benar, sehingga
\[\tf{\land}(\False,\Undef)=
\tf{\land}(\Undef,\False)=\False.\]
Nilai-nilai lainnya, yaitu ketika argumennya merupakan nilai kebenaran yang
sudah ditentukan, $\True$ atau~$\False$, sama seperti dalam logika klasik.

Bagi implikasi, keadaannya sedikit lebih rumit. Andaikan $q$ merupakan
pernyataan kontingen masa depan. Jika $p$ salah, maka $p\lif q$ akan benar
bagaimanapun hasil~$q$, sehingga kita harus menetapkan
$\tf{\lif}(\False,\Undef)=\True$. Jika $p$ benar, maka $q\lif p$ akan benar
bagaimanapun hasil~$q$, sehingga $\tf{\lif}(\Undef,\True)=\True$. Jika $p$
benar, maka $p\lif q$ mungkin menjadi benar atau salah, sehingga
$\tf{\lif}(\True,\Undef)=\Undef$. Serupa dengan itu, jika $p$ salah, maka
$q\lif p$ mungkin menjadi benar atau salah, sehingga
$\tf{\lif}(\Undef,\False)=\Undef$. Tersisa kasus ketika $p$ dan~$q$
keduanya merupakan kontingen masa depan. Berdasarkan motivasinya, dalam kasus
ini kita seharusnya menetapkan~$\Undef$. Namun, hal itu akan membuat
$!A\lif !A$ \emph{bukan} suatu tautologi. \L ukasiewicz tidak keberatan
melepaskan $!A\lor\lnot !A$ dan $\lnot(!A\land\lnot !A)$, tetapi enggan
melepaskan $!A\lif !A$. Karena itu, ia menetapkan
$\tf{\lif}(\Undef,\Undef)=\True$.

\begin{defn}\ollabel{def:lukasiewicz}
Logika \L ukasiewicz bernilai tiga didefinisikan dengan menggunakan matriks
berikut:
\begin{enumerate}
  \item Bahasa proposisional baku~$\Lang L_0$ dengan $\lnot$, $\land$,
  $\lor$, $\lif$.
  \item Himpunan nilai kebenaran~$V=\{\True,\Undef,\False\}$.
  \item $\True$ merupakan satu-satunya nilai yang ditetapkan, yakni
  $V^+=\{\True\}$.
  \item Fungsi kebenaran diberikan oleh tabel-tabel berikut:
  \begin{center}
    \begin{tabular}{c|c}
      $\tf{\lnot}$ & \\
      \hline
      $\True$ & $\False$ \\
      $\Undef$ & $\Undef$ \\
      $\False$ & $\True$
    \end{tabular}
    \quad
    \begin{tabular}{c|ccc}
      $\tf{\land}[\LogLuk[3]]$ & $\True$ & $\Undef$ & $\False$ \\
      \hline
      $\True$ & $\True$ & $\Undef$ & $\False$ \\
      $\Undef$ & $\Undef$ & $\Undef$ & $\False$\\
      $\False$ & $\False$ & $\False$ & $\False$
    \end{tabular}
    \\[2ex]
    \begin{tabular}{c|ccc}
      $\tf{\lor}[\LogLuk[3]]$ & $\True$ & $\Undef$ & $\False$ \\
      \hline
      $\True$ & $\True$ & $\True$ & $\True$ \\
      $\Undef$ & $\True$ & $\Undef$ & $\Undef$ \\
      $\False$ & $\True$ & $\Undef$ & $\False$
    \end{tabular}
    \quad
    \begin{tabular}{c|ccc}
      $\tf{\lif}[\LogLuk[3]]$ & $\True$ & $\Undef$ & $\False$ \\
      \hline
      $\True$ & $\True$ & $\Undef$ & $\False$ \\
      $\Undef$ & $\True$ & $\True$ & $\Undef$  \\
      $\False$ & $\True$ & $\True$ & $\True$
    \end{tabular}
  \end{center}
\end{enumerate}
\end{defn}

Mudah dilihat bahwa setiap !!{formula}~$!A$ yang hanya memuat $\lnot$,
$\land$, dan~$\lor$ akan bernilai~$\Undef$ jika semua
!!{propositional variable}s di dalamnya diberi nilai~$\Undef$. Jadi, misalnya,
tautologi klasik
$p\lor\lnot p$ dan $\lnot(p\land\lnot p)$ bukan tautologi
dalam~$\LogLuk[3]$, karena $\pValue{v}(!A)=\Undef$ setiap kali
$\pAssign v(p)=\Undef$.

Pada !!{valuation}s dengan $\pAssign v(p)=\True$ atau~$\False$,
$\pValue v(!A)$ sama dengan nilai kebenaran klasiknya.

\begin{prop}
  Jika $\pAssign v(p)\in\{\True,\False\}$ bagi semua $p$ dalam~$!A$, maka
  $\pValue v(!A)[\LogLuk[3]]=\pValue v(!A)[\LogCL]$.
\end{prop}

\begin{prob}\label{mvl:thr:luk:prob:luk-iff} Andaikan kita mendefinisikan
  $\pValue v(!A\liff !B)=\pValue v((!A\lif !B)\land(!B\lif
  !A))$ dalam~$\LogLuk[3]$. Tabel kebenaran apakah yang akan dimiliki
  oleh~$\liff$?
\end{prob}

Banyak tautologi klasik \emph{juga} merupakan tautologi dalam~\LogLuk[3],
misalnya $\lnot p\lif(p\lif q)$. Sama seperti dalam logika klasik, kita dapat
memverifikasinya dengan tabel kebenaran:
\begin{center}
\begin{tabular}{cc|cccccc}
  $p$ & $q$ & $\lnot$ & $p$ & $\lif$ & $\smash{(}p$ & $\lif$ & $q\smash{)}$ \\ \hline
  \True & \True & \False & \True & \True & \True & \True & \True \\
  \True & \Undef & \False & \True & \True & \True & \Undef & \Undef \\
  \True & \False & \False & \True & \True & \True & \False & \False \\
  \Undef & \True & \Undef & \Undef & \True & \Undef & \True & \True \\
  \Undef & \Undef & \Undef & \Undef & \True & \Undef & \True & \Undef \\
  \Undef & \False & \Undef & \Undef & \True & \Undef & \Undef & \False \\
  \False & \True & \True & \False & \True & \False & \True & \True \\
  \False & \Undef & \True & \False & \True & \False & \True & \Undef \\
  \False & \False & \True & \False & \True & \False & \True & \False \\
\end{tabular}
\end{center}

\begin{prob}
  Tunjukkan bahwa formula-formula berikut merupakan tautologi
  dalam~\LogLuk[3]:
  \begin{enumerate}
    \item $p\lif(q\lif p)$
    \item\label{mvl:thr:luk:prob:luk-taut-2} $\lnot(p\land q)\liff(\lnot p\lor\lnot q)$
    \item\label{mvl:thr:luk:prob:luk-taut-3} $\lnot(p\lor q)\liff(\lnot p\land\lnot q)$
  \end{enumerate}
  (Dalam \olref[mvl][thr][luk]{prob:luk-taut-2} dan
  \olref[mvl][thr][luk]{prob:luk-taut-3}, gunakan $!A\liff !B$ sebagai
  singkatan bagi $(!A\lif !B)\land(!B\lif !A)$, atau rujuk jawaban Anda
  atas \olref[mvl][thr][luk]{prob:luk-iff}.)
\end{prob}

\begin{prob}
  Tunjukkan bahwa tautologi-tautologi klasik berikut bukan tautologi
  dalam~\LogLuk[3]:
  \begin{enumerate}
    \item $(\lnot p\land p)\lif q$
    \item $((p\lif q)\lif p)\lif p$
    \item $(p\lif(p\lif q))\lif(p\lif q)$
  \end{enumerate}
\end{prob}

Karena itu, mungkin orang akan menduga bahwa meskipun tidak semua tautologi
klasik merupakan tautologi dalam~$\LogLuk[3]$, setiap tautologi klasik
setidaknya harus bernilai~$\True$ atau~$\Undef$ pada setiap !!{valuation}.
Namun, dugaan ini salah. Suatu contoh penyangkal diberikan oleh
\[
  \lnot(p\lif\lnot p)\lor\lnot(\lnot p\lif p)
\]
yang bernilai~$\False$ jika $p$ bernilai~$\Undef$.

\begin{prob}
  Manakah di antara relasi-relasi berikut yang berlaku dalam logika
  \L ukasiewicz? Berikan tabel kebenaran bagi masing-masing relasi.
  \begin{enumerate}
    \item $p,p\lif q\Entails q$
    \item $\lnot\lnot p\Entails p$
    \item $p\land q\Entails p$
    \item $p\Entails p\land p$
    \item $p\Entails p\lor q$
  \end{enumerate}
\end{prob}

\L ukasiewicz berharap dapat membangun suatu logika kemungkinan berdasarkan
sistem bernilai tiganya dengan memperkenalkan konektif uner
$\Diamond !A$ (untuk ``$!A$ mungkin'') dan konektif yang bersesuaian
$\Box !A$ (untuk ``$!A$ niscaya''):
\begin{center}
  \begin{tabular}{c|c}
    $\tf{\Diamond}$ & \\
    \hline
    $\True$ & $\True$ \\
    $\Undef$ & $\True$ \\
    $\False$ & $\False$
  \end{tabular}
  \quad
  \begin{tabular}{c|c}
    $\tf{\Box}$ & \\
    \hline
    $\True$ & $\True$ \\
    $\Undef$ & $\False$ \\
    $\False$ & $\False$
  \end{tabular}
\end{center}
Dengan kata lain, $p$ mungkin jika dan hanya jika nilainya belum ditetapkan
sebagai salah; dan $p$ niscaya jika dan hanya jika nilainya sudah ditetapkan
sebagai benar.

\begin{prob}
  Tunjukkan bahwa $\Box p\liff\lnot\Diamond\lnot p$ dan
  $\Diamond p\liff\lnot\Box\lnot p$ merupakan tautologi dalam~\LogLuk[3]
  yang diperluas dengan tabel kebenaran bagi $\Box$ dan~$\Diamond$.
\end{prob}

Akan tetapi, kekurangan usulan logika modal ini segera tampak: bagaimanapun
keadaan akhirnya, $p\land\lnot p$ tidak pernah mungkin menjadi benar. Jadi,
meskipun nilainya sekarang belum ditetapkan (dan karena itu bernilai tak tentu),
formula tersebut seharusnya dianggap mustahil; dengan kata lain,
$\lnot\Diamond(p\land\lnot p)$ seharusnya merupakan tautologi. Namun, jika
$\pAssign v(p)=\Undef$, maka
$\pValue v(\lnot\Diamond(p\land\lnot p))=\False$. Meskipun \L ukasiewicz
benar bahwa dua nilai kebenaran tidak cukup untuk menampung pembedaan modal
seperti kemungkinan dan keniscayaan, memperkenalkan nilai kebenaran ketiga
juga belum cukup.

\end{document}
